\chapter{Thực nghiệm, đánh giá và thảo luận}
\label{chap:evaluation}

\section{Thiết lập thực nghiệm}

\subsection{Môi trường phần cứng}

% TODO: Điền thông tin GPU, RAM, storage thực tế
\begin{table}[h]
\centering
\renewcommand{\arraystretch}{1.3}
\caption{Cấu hình phần cứng cho huấn luyện và inference}
\label{tab:hardware}
\resizebox{0.6\textwidth}{!}{%
\begin{tabular}{|c|c|}
\hline
\textbf{Thành phần} & \textbf{Cấu hình} \\
\hline
GPU & \ldots (VD: NVIDIA A100 40GB) \\
\hline
RAM & \ldots GB \\
\hline
Storage & \ldots GB SSD \\
\hline
OS & Ubuntu 22.04 LTS \\
\hline
\end{tabular}%
}
\end{table}

\subsection{Môi trường phần mềm}

\begin{itemize}
    \item Python 3.10+, CUDA 12.x
    \item LLaMA-Factory (SFT framework)
    \item vLLM (inference engine)
    \item FastAPI, PostgreSQL, Docker, Judge0
\end{itemize}

\section{Kết quả đánh giá DebugEval}

\subsection{Hiệu năng Gemma-4 (Base vs. Fine-Tuned)}

\begin{table}[h]
\centering
\renewcommand{\arraystretch}{1.3}
\caption{Kết quả Gemma-4-E4B-IT trên DebugEval Benchmark}
\label{tab:gemma_results}
\resizebox{\textwidth}{!}{%
\begin{tabular}{|c|c|c|c|c|}
\hline
\textbf{Mô hình} & \textbf{Task 1 (\%)} & \textbf{Task 2 (\%)} & \textbf{Task 3 (\%)} & \textbf{Task 4 (\%)} \\
\hline
Gemma4 Base & 79.24 & 55.43 & 8.94 & N/A* \\
\hline
Gemma4 + LoRA & 73.01 & 44.61 & \textbf{41.30} ($\uparrow$×4.6) & 86.65 \\
\hline
\end{tabular}%
}
\end{table}

\textit{*Gemma4 Base Task 4 bị gián đoạn và không cho ra kết quả cuối cùng.}

\textbf{Nhận xét:} Gemma-4 sau LoRA cho thấy \textbf{sự đánh đổi chuyên biệt hóa} (specialization trade-off): Task 3 (Code Repair) cải thiện vượt trội 4.6 lần, nhưng Task 1, 2 (comprehension) giảm nhẹ 6--11 điểm phần trăm. Điều này cho thấy việc chỉ tinh chỉnh $q\_proj$, $v\_proj$ có thể làm mô hình thiên về kỹ năng generative thay vì discriminative.

\subsection{Hiệu năng Qwen2.5-Coder (Base vs. Fine-Tuned)}

\begin{table}[h]
\centering
\renewcommand{\arraystretch}{1.3}
\caption{Kết quả Qwen2.5-Coder-3B-Instruct trên DebugEval Benchmark}
\label{tab:qwen_results}
\resizebox{\textwidth}{!}{%
\begin{tabular}{|c|c|c|c|c|}
\hline
\textbf{Mô hình} & \textbf{Task 1 (\%)} & \textbf{Task 2 (\%)} & \textbf{Task 3 (\%)} & \textbf{Task 4 (\%)} \\
\hline
Qwen Base & 44.29 & 25.43 & 34.30 & 83.38 \\
\hline
Qwen + LoRA & \textbf{71.45} ($\uparrow$+27) & \textbf{43.79} ($\uparrow$+18) & \textbf{44.69} ($\uparrow$+10) & 83.75 \\
\hline
\end{tabular}%
}
\end{table}

\textbf{Nhận xét:} Qwen Coder + LoRA đạt \textbf{cải thiện đồng đều trên cả 4 task} mà không có suy giảm, là cấu hình cân bằng nhất.

\subsection{So sánh chéo giữa các mô hình}

\textbf{Phát hiện quan trọng:} Sau fine-tuning, cả hai mô hình (dù khởi đầu với base capabilities rất khác nhau) \textbf{hội tụ về hiệu năng tương đồng} (chênh lệch trong $\pm$3 điểm phần trăm). Điều này gợi ý rằng chất lượng dữ liệu fine-tuning có ảnh hưởng lớn hơn kích thước mô hình trong phạm vi benchmark này.

\section{Kết quả đánh giá Socratic Tutoring}

\subsection{Đánh giá bằng Heuristic Metrics}

\begin{table}[h]
\centering
\renewcommand{\arraystretch}{1.3}
\caption{Kết quả đánh giá Socratic Tutor (Qwen2.5-3B + LoRA)}
\label{tab:socratic_results}
\resizebox{0.7\textwidth}{!}{%
\begin{tabular}{|c|c|}
\hline
\textbf{Metric} & \textbf{Kết quả} \\
\hline
Question Rate (tỷ lệ đặt câu hỏi dẫn dắt) & \textbf{96.1\%} \\
\hline
Answer Concealment (tỷ lệ che giấu đáp án) & \textbf{100.0\%} \\
\hline
\end{tabular}%
}
\end{table}

\subsection{Đánh giá bằng LLM-as-Judge (Gemma-4)}

\begin{table}[h]
\centering
\renewcommand{\arraystretch}{1.3}
\caption{Điểm đánh giá từ Gemma-4 Judge (thang 5)}
\label{tab:llm_judge}
\resizebox{0.5\textwidth}{!}{%
\begin{tabular}{|c|c|}
\hline
\textbf{Tiêu chí} & \textbf{Điểm trung bình} \\
\hline
On-Topic Relevance & \textbf{4.93}/5 \\
\hline
Helpfulness & \textbf{4.49}/5 \\
\hline
\end{tabular}%
}
\end{table}

\textbf{Nhận xét:} Mô hình tinh chỉnh đạt tỷ lệ che giấu đáp án tuyệt đối (100\%), vượt mục tiêu NFR.5 ($>$98\%). Tỷ lệ đặt câu hỏi 96.1\% cho thấy mô hình duy trì phong cách Socratic ổn định.

\section{Đánh giá hệ thống tích hợp}

\subsection{Kiểm thử chức năng}

% TODO: Bổ sung kết quả test thực tế
\begin{table}[h]
\centering
\renewcommand{\arraystretch}{1.3}
\caption{Kết quả kiểm thử các chức năng chính}
\label{tab:func_test}
\resizebox{\textwidth}{!}{%
\begin{tabular}{|c|c|c|}
\hline
\textbf{Chức năng} & \textbf{Mô tả test} & \textbf{Kết quả} \\
\hline
FR.1 Context Capture & Bôi đen code, gửi chat & PASS \\
\hline
FR.2 Socratic Chat & Chat nhận phản hồi Socratic & PASS \\
\hline
FR.5 Authentication & Đăng nhập/đăng xuất trên Extension & PASS \\
\hline
FR.6 Assignment TreeView & Hiển thị danh sách bài tập & PASS \\
\hline
FR.8 Submit \& Summarize & Nộp bài, nhận phản hồi Judge + AI & PASS \\
\hline
FR.9 Course Management & CRUD lớp/bài tập trên Web & PASS \\
\hline
FR.14 Judge Server & Chạy code trong sandbox & PASS \\
\hline
\end{tabular}%
}
\end{table}

\subsection{Kiểm thử hiệu năng}

\begin{table}[h]
\centering
\renewcommand{\arraystretch}{1.3}
\caption{Kết quả kiểm thử yêu cầu phi chức năng}
\label{tab:nfr_test}
\resizebox{\textwidth}{!}{%
\begin{tabular}{|c|c|c|c|}
\hline
\textbf{Chỉ tiêu} & \textbf{Mục tiêu} & \textbf{Thực tế} & \textbf{Đạt?} \\
\hline
NFR.1 Zero-friction Access & $<$ 10s & \ldots s & \ldots \\
\hline
NFR.2 AI Response Time & $<$ 3s & \ldots s & \ldots \\
\hline
NFR.3 Judge + AI Response & $<$ 7s & \ldots s & \ldots \\
\hline
NFR.5 Anti-Solution Rate & $>$ 98\% & 100\% & PASS \\
\hline
\end{tabular}%
}
\end{table}

\section{Thảo luận}

\subsection{Ưu điểm}

\begin{enumerate}
    \item Hệ thống đạt tỷ lệ che giấu đáp án 100\%, đảm bảo nguyên tắc Socratic.
    \item Kỹ thuật LoRA giúp tinh chỉnh mô hình hiệu quả với chi phí phần cứng thấp.
    \item Kiến trúc LangGraph cho phép quản lý luồng AI tường minh, dễ debug.
    \item Tích hợp trọn vẹn vào VS Code giúp sinh viên không cần rời IDE.
\end{enumerate}

\subsection{Hạn chế}

\begin{enumerate}
    \item Gemma-4 + LoRA có hiện tượng trade-off giữa các task, cần nghiên cứu thêm multi-task fine-tuning.
    \item Bộ dữ liệu Socratic còn nhỏ (552 mẫu train), có thể giới hạn khả năng tổng quát hóa.
    \item Chưa thực hiện user study quy mô lớn với sinh viên thực tế.
    \item Hệ thống phụ thuộc vào GPU cho inference mô hình tinh chỉnh.
\end{enumerate}

\subsection{Phân tích lỗi}

% TODO: Bổ sung phân tích các trường hợp lỗi cụ thể

Một số trường hợp mô hình Socratic không đạt yêu cầu:
\begin{itemize}
    \item 3.9\% phản hồi không chứa câu hỏi dẫn dắt, thường xảy ra khi lỗi quá đơn giản (typo, thiếu dấu chấm phẩy).
    \item Với lỗi phức tạp liên quan đến thuật toán nâng cao, mô hình 3B đôi khi phân tích root cause chưa chính xác.
\end{itemize}
